\paragraph{Worked semantic calculation.}
At a normal world \(n\in N\), the displayed base condition reduces implication
to global inclusion of truth sets:
\[
\begin{aligned}
  \mathcal M,n\entails A\to B
  &\Longleftrightarrow
    \forall u,v\,
      ((u=v\land\mathcal M,u\entails A)
        \Longrightarrow\mathcal M,v\entails B)\\
  &\Longleftrightarrow
    \forall u\in W\,
      (\mathcal M,u\entails A
       \Longrightarrow\mathcal M,u\entails B).
\end{aligned}
\]
Thus implication at \(n\) is not evaluated merely from the truth values of
\(A\) and \(B\) at \(n\).  At a nonnormal world, the ternary relation can
instead test an \(A\)-world and a possibly different \(B\)-world.  The
four-world countermodel below exploits exactly that difference.

For negation, the star operation supplies the comparison point.  Let
\(W=\set{w,w^*}\), let \(N=\set w\), let star swap the two worlds, put
\[
  R=\set{\tuple{w,x,x}:x\in W},
  \qquad
  V(p)=\set w,
\]
and let \(\mathcal M=(W,N,R,{*},V)\), with the values of other
propositional variables arbitrary.
Then \(\mathcal M,w\entails\neg p\) holds while
\(\mathcal M,w^*\notentails\neg p\).  This is a direct calculation from the
semantic clause, not an assertion that every relevance logic validates the
same additional negation principles.

